\documentclass[UTF8]{ctexart}
\usepackage{amsmath,amssymb}
\begin{document}
\section*{公共交通货运战略规划：路径流模型}

本文件仅采用原选定完整模型的路径流表述。站点流和线路状态流是另选表述，
不与路径流式混成一个模型。符号补齐依据原论文第5、6、8页。
$F,L,S,K$分别为车队、线路、站点、商品；$f(\ell)$为线路所属车队，
$L_f=\{\ell:f(\ell)=f\}$，$L_s=\{\ell:s\in S_\ell\}$。
每条有向线路的$S_\ell^{in}$在路径上先于$S_\ell^{out}$，两者不交；
$S_{\ell k}^{in}=S_\ell^{in}\cap S_k^{in}$，
$S_{\ell k}^{out}=S_\ell^{out}\cap S_k^{out}$。
定义$\mathcal W=\{(k,\ell,s,t):s\in S_{\ell k}^{in},t\in S_{\ell k}^{out}\}$。
变量$x_s,y_\ell,z_f$为0--1，$w_{k\ell st}\in\mathbb Z_{\ge0}$仅对$\mathcal W$定义，
求和域之外按0处理。
\begin{align*}
\max\quad &\sum_{(k,\ell,s,t)\in\mathcal W}w_{k\ell st}\\
\text{s.t.}\quad&
\sum_s V_sx_s+\sum_\ell V_\ell y_\ell+\sum_f V_fz_f\le B,\\
&\sum_s x_s\le\mathrm{maxstops},\quad
\sum_\ell y_\ell\le\mathrm{maxlines},\quad\sum_fz_f\le\mathrm{maxfleets},\\
&y_\ell\le z_{f(\ell)}&&(\ell\in L),\\
&x_s\le\sum_{\ell\in L_s}y_\ell&&(s\in S),\\
&y_\ell\le\sum_{s\in S_\ell^{in}}x_s,\quad
y_\ell\le\sum_{s\in S_\ell^{out}}x_s&&(\ell\in L),\\
&z_f\le\sum_{\ell\in L_f}y_\ell&&(f\in F),\\
&\sum_{(\ell,s,t):(k,\ell,s,t)\in\mathcal W}w_{k\ell st}\le P_k&&(k\in K),\\
&\sum_{(k,s,t):(k,\ell,s,t)\in\mathcal W}w_{k\ell st}\le C_\ell y_\ell&&(\ell\in L),\\
&\sum_{(k,\ell,u,v)\in\mathcal W}
(\mathbf1_{\{u=s\}}+\mathbf1_{\{v=s\}})w_{k\ell uv}\le C_sx_s&&(s\in S).
\end{align*}
补全实例保留线路站点顺序、需求、容量和预算；对缺少的线路上下货角色，
将每条线路前半站点指定为上货、后半站点为下货，作为明确的运营数据补全。
站点角色由上述数据定义。
禁止同站上/下货或逆序运输；最大覆盖模型允许部分商品不被覆盖。

\end{document}
